% Include after the positive-angular main-Gram construction.
% Requires amsmath, amssymb, amsthm and mathrsfs; theorem is already defined.
\section{Fixed testing functions for the complete positive-angular parent}
Retain the preceding source/testing rectangle and $\lambda>0$.
Compose the testing functions with the angle gauge already defined, then
multiply all four state coordinates by $\operatorname{sgn}_+(W)$,
where $\operatorname{sgn}_+(0)=1$. Thus $W\ge0$.
Put $Z_+=\max(Z,0)$ and $Z_-=\max(-Z,0)$. The common state sign
preserves the quadratic parent. Every testing function below is pulled
back through this one fixed map.

\subsection{The exact whole-parent decomposition}

All symbols $g,G,N_x,N_y,R,\Gamma,A_\theta,T_\theta,J_C,J_\ell$
have their explicit definitions in \texttt{fixed\_box\_c1.tex}.
Set $F_A=u(Z-a_0wW)$, $X_e=X+ewW$, and
\[
 a_C=cd_G/S,\quad q_C^{(0)}=\tfrac12-a_C(2/\lambda+k),\quad
 p_C=q_C^{(0)}-a_C\delta.
\]
The proved bounds give $0\le a_C\le6$ and $p_C>1/4$.
The circuit remainder satisfies
\[
 \mathscr C=\sum_{i=1}^4d_i^2+\mathscr C_+,
\]
and the angular difference is
\[
 \Delta_A=(2/\lambda+u+k)\mathscr C+u^2R^2
       +\{2a_0u+(u+k)/L_0\}\Gamma^2-A_\theta^2.
\]
The complete parent is exactly
\begin{align*}
 F_w={}&g^{\mathsf T}Gg+\mathcal E_+
 +(2+4m)w^2N_x^2
 +\{4a+2z(2C+1)-1\}J_C^2+(4\ell-1)J_\ell^2\\
 &+\tfrac12L_0F_A^2
 +{K_A+c-1\over2}(A_\theta^2+T_\theta^2)
 +{K_A+c\over2}(A_\theta^2-T_\theta^2)\\
 &+{cd_G\over S}\Delta_A
 +\{p_C+a_C(\delta-u)\}\mathscr C_+,
\end{align*}
where
\[
\begin{split}
 \mathcal E_+={}&4\beta(1+w^2)N_x^2+2(2D+1)(uX_e)^2
                 +4bpzu^2\Gamma^2\\
 &+4a(u^2+w^2+u^2w^2)J_C^2+4\ell u^2J_\ell^2
                                             +4Czw^2J_C^2.
\end{split}
\]
After the expansions below, every square coefficient outside $G$
is nonnegative and source-only, except the explicitly separated
testing factor $\delta-u$, which is nonnegative on the fixed
testing rectangle. This equality
is the full-state identity, not its restriction to a stationary
state or a determinant. No inverse of $D$ or $\beta$ occurs.

\subsection{Ordinary retained squares: seven additional functions}

Let $\mathcal V$ be the twelve pre-Gram functions
\[
 \{X,Y,Z,W,uZ,wZ,uW,uwW,u^2W,
                     u^{3/2}W,u^{1/2}W,\sqrt{uwWZ_+}\}.
\]
All ordinary retained roots are source-linear combinations of
$\mathcal V$ and the seven functions
\[
                   wX,\ uwZ,\ uX,\ u^2wW,\ uY,\ wY,\ uwY.
\]
Here are all seventeen roots; every displayed square root of a
coefficient is source-only:
\[
\begin{array}{c|l}
\text{family}&\text{square roots}\\\hline
\mathcal E_+&2\sqrt\beta N_x,\quad2\sqrt\beta\,wN_x,\quad
 \sqrt{2(2D+1)}\,uX_e,\quad2\sqrt{bpz}\,u\Gamma\\
 &2\sqrt a\,uJ_C,\quad2\sqrt a\,wJ_C,\quad2\sqrt a\,uwJ_C,\quad
 2\sqrt\ell\,uJ_\ell,\quad2\sqrt{Cz}\,wJ_C\\
N_x\text{ excess}&\sqrt{2+4m}\,wN_x\\
\text{radial excess}&\sqrt{4a+2z(2C+1)-1}\,J_C,\quad
                              \sqrt{4\ell-1}\,J_\ell\\
\text{averaged square}&\sqrt{L_0/2}\,F_A\\
\text{angular excess}&\sqrt{(K_A+c-1)/2}\,A_\theta,\quad
                       \sqrt{(K_A+c-1)/2}\,T_\theta\\
\text{angular difference}&\sqrt{(K_A+c)/2}\,qu(u-k)W,\\
 &\sqrt{(K_A+c)/2}\,q\sqrt{2k}\sqrt u(u-k)W.
\end{array}
\]
For example $J_C=Y/q-bquW$ and
$J_\ell=Y/q-wZ/q+QquW$, so multiplying by $u,w,uw$ adds
exactly the stated functions. All other rows follow by substitution.
The last two squares use
$A_\theta^2-T_\theta^2=z(u-k)^2W^2(u^2+2uk)$.

\subsection{The angular difference: nineteen functions}

The following calculation clears the intermediate state ratio
completely. Define positive source constants
\[
 A_u=1/k,\quad A_v=1/(\sqrt\sigma q),\quad
 A_z=\sqrt\sigma/(kq),\quad M_0=k^2q,\quad a_*=\lambda k,
\]
and write $r_* =\sqrt{wW|Z|}$. Eight cleared roots are
\[
\begin{split}
 f_1={}&M_0(A_uA_vA_zuw|Z|-W),\\
 f_2={}&\sqrt2 M_0(A_vA_zw|Z|-A_uuW),\\
 f_3={}&M_0\{(\sqrt{a_*}+1/\sqrt{a_*})A_uA_vuwW
                   -\sqrt{a_*}A_uA_zu|Z|-A_z|Z|/\sqrt{a_*}\},\\
 f_4={}&\sqrt{A_u}\sqrt u\,f_3,\\
 f_5={}&\sqrt{a_*}M_0(W-A_vA_zw|Z|),\qquad
 f_6=\sqrt{A_u}\sqrt u\,f_5,\\
 f_7={}&\sqrt{2a_*A_vA_z}M_0(A_uu-1)r_*,\qquad
 f_8=\sqrt{A_u}\sqrt u\,f_7.
\end{split}
\]
For $Z\ge0$, $\Delta_A=\sum_{i=1}^8f_i^2$. To verify this,
put $U=A_uu$, $V=A_vw$, $Z_*=A_z|Z|/W$, $M=M_0W$;
the roots are $M$ times the eight roots obtained by expanding
every $U+1$ factor in the identity
\begin{align*}
 &[2+a_*(U+1)]\Psi+U^2Q_*^2
       +[2U+(U+1)/a_*]R_*^2-T_*^2\\
 &=(T_*+UQ_*)^2+2Q_*^2
 +(U+1)(\sqrt{a_*}U(V-Z_*)-R_*/\sqrt{a_*})^2\\
 &\quad+a_*(U+1)(1-VZ_*)^2
                         +2a_*(U+1)VZ_*(U-1)^2,
\end{align*}
where $\Psi=1+U^2V^2+U^2Z_*^2+V^2Z_*^2-4UVZ_*$,
$Q_*=VZ_*-U$, $R_*=Z_*-UV$, and $T_*=U^2-1$.
It is a polynomial identity by multiplication.

For $Z<0$, add three squares. Put
\[
 C_2=4,\quad C_1=8a_*+8+4/a_*,\quad C_0=16+8a_*+4/a_*.
\]
Their roots are
\[
 M_0\sqrt{A_uA_vA_z}\,\boldsymbol1_{Z<0}
 \bigl(\sqrt{C_2}A_u u^{3/2}r_*,\quad
       \sqrt{C_1A_u}\,u r_*,\quad
       \sqrt{C_0}\sqrt u\,r_*\bigr).
\]
Indeed changing $Z_*$ to $-Z_*$ increases the left side by
$M^2UV|Z_*|(C_2U^2+C_1U+C_0)$. This also proves the exact
negative-state-sector coefficients. Every displayed source
factor is a real closed field square root in the source factor.

The resulting nineteen testing functions, including overlaps
with $\mathcal V$, are
\[
\begin{split}
 \mathcal A=\{&W,uw|Z|,w|Z|,uW,uwW,u|Z|,|Z|,
 u^{3/2}wW,u^{3/2}|Z|,\sqrt u|Z|,\sqrt uW,\sqrt u w|Z|,\\
 &ur_*,r_*,u^{3/2}r_*,\sqrt u r_*,
 \boldsymbol1_{Z<0}u^{3/2}r_*,
 \boldsymbol1_{Z<0}ur_*,\boldsymbol1_{Z<0}\sqrt u r_*\}.
\end{split}
\]
Multiply the eleven root formulas above by the pure-source factor
$\sqrt{cd_G/S}$ for their contribution to $F_w$.
Every cleared formula is defined at $W=0$ and $Z=0$. After
substituting squares for $u,w,W,|Z|$, the identities are polynomial;
therefore they include those literal boundaries, without a limit
or a state-coordinate inverse in the final formulas.

\subsection{The circuit remainder: fourteen functions}

Set the positive source constants
\[
 s_1=\sqrt{\lambda\sigma},\quad s_2=\sqrt\rho,\quad
 s_3=\sqrt\lambda\,xq,\quad s_4=\sqrt{\lambda x}/q,
 \qquad g_0=\sqrt{\lambda x}.
\]
Define seven testing functions
\[
\begin{array}{ll}
 t_1=u^{3/4}w^{1/4}W^{1/4}Z_+^{3/4},&
 t_2=u^{3/4}w^{3/4}W^{3/4}Z_+^{1/4},\\
 t_3=u^{1/4}w^{1/4}W^{3/4}Z_+^{1/4},&
 t_4=u^{1/4}w^{3/4}W^{1/4}Z_+^{3/4},\\
 t_5=u^{3/4}w^{1/2}W^{1/2}Z_+^{1/2},&
 t_6=u^{1/4}w^{1/2}W^{1/2}Z_+^{1/2},\\
 t_7=\sqrt{uwWZ_-}.&
\end{array}
\]
An explicit full-sector SOS is
\[
 \mathscr C_+=r_1^2+r_2^2+r_3^2+r_4^2,
\]
\[
\begin{gathered}
 r_1=\sqrt{2g_0}(\sqrt{s_1}t_1-\sqrt{s_2}t_2),\quad
 r_2=\sqrt{2g_0}(\sqrt{s_3}t_3-\sqrt{s_4}t_4),\\
 r_3=2\sqrt{g_0}\{(s_1s_2)^{1/4}t_5-(s_3s_4)^{1/4}t_6\},
 \qquad r_4=2\sqrt{\lambda x}\,t_7.
\end{gathered}
\]
The first three roots vanish for $Z\le0$ and the last vanishes
for $Z\ge0$. They also vanish wherever $WZ=0$. On $Z>0$ they
are precisely the separated two-pairing four-monomial identity;
on $Z<0$ the remainder is $4\lambda xuwWZ_-$.

There is no point-dependent scalar reserve at this step:
\[
 (p_C+a_C(\delta-u))\mathscr C_+
      =\sum_{j=1}^4(\sqrt{p_C}r_j)^2
        +\sum_{j=1}^4(\sqrt{a_C}\sqrt{\delta-u}\,r_j)^2.
\]
Thus the fixed circuit feature list is exactly the seven $t_i$
and their seven products with $\sqrt{\delta-u}$.

\subsection{The evaluated union and machine-readable data}

For the main Gram take the seven nonconstant monomials
\[
 \mathcal B=\{w,w^2,\sqrt u,u,uw,u^2,u^2w\},\qquad
 \gamma_7={1\over7\,2^{38}}.
\]
The main list is $\mathcal V$ times the fifteen fixed factors
$1,\sqrt{\gamma_7+b},\sqrt{\gamma_7-b}$ for $b\in\mathcal B$,
giving $180$ functions. Its $144$ fixed directions are $e_i$,
$e_i+e_j$, and $e_i-e_j$, with the exact scalar-channel formulas
of \texttt{fixed\_box\_c1.tex}, using $7A_*\gamma_7=c_0$.

\begin{theorem}
On the stated independent rectangle, the full real-state parent
$F_w$ is a sum of squares of source-linear combinations of one
fixed list of $220$ labelled testing functions: the $180$ main
functions, the seven ordinary functions, the nineteen angular
functions, and the fourteen circuit functions. Repeated functions
are harmless; the real dimension of their span is at most $220$.
There is an explicit construction using at most $2196$ square
roots: $144\cdot15$ main roots, seventeen ordinary roots, eleven
angular roots, and eight circuit roots. The same list, composed
with the fixed gauges above, includes all real state-sign sectors,
both testing axes, and the source faces $D=0$ and $\beta=0$.
\end{theorem}

This theorem follows from the displayed exact decomposition and
the separated coefficient formulas. It needs no Gram factorization
with mixed radicals. Without using the physical reality reduction,
apply the theorem separately to the real and imaginary state
amplitudes; their union has at most $440$ functions and $4392$
roots. The parent is real quadratic in the state, so those two
contributions sum to its complex-state value.

The adjacent \texttt{c1\_retained\_union\_checks.py} writes all
$220$ labelled features, all $144$ main directions, the main
coefficient matrix and root formulas to
\texttt{c1\_positive\_lambda\_local220.json}. It verifies the two
normal identities, the angular identity and signed correction,
the circuit identity, the scalar circuit split, the whole-parent
assembly, the main matrix support, the fixed-direction identity,
and the feature/root counts by exact symbolic arithmetic.
This is an evaluated local testing dictionary, not a global
full-set affine pencil: the other branches and global cover
selection remain outside its scope.
